\section{Introduction}

% Burden of stroke

Stroke remains one of the top three global causes of death and disability, accounting for 10.7\% of all deaths globally and 5.7\% of lost disability-adjusted life-years (DALYs) \cite{feigin_global_2021}. Despite reductions in age-standardised rates of stroke, ageing populations are driving an increase in the absolute number of strokes \cite{feigin_global_2021}. Across Europe, in 2017, stroke was found to cost healthcare systems \texteuro 27 billion, or 1.7\% of health expenditure \cite{luengo-fernandez_economic_2020}. The burden is increasing; it has been estimated that the number of stroke survivors aged 45 and over in the UK will more than double between 2015 and 2035 \cite{king_future_2020}.

% IVT

Intravenous thrombolysis (IVT) with alteplase was established by randomised trials showing benefit when given within 3 hours of onset \cite{ninds1995}, a window later extended to 4.5 hours \cite{hacke2008}. The individual patient data meta-analysis of Emberson et al., pooling 6{,}756 patients from nine trials, quantified how strongly that benefit depends on time \cite{emberson_effect_2014}. The odds of an excellent outcome (a modified Rankin Scale [mRS] score of 0--1 at 3--6 months) were increased by 75\% when treatment was given within 3 hours (OR 1.75, 95\% CI 1.35--2.27) and by 26\% between 3 and 4.5 hours (OR 1.26, 95\% CI 1.05--1.51), with no significant benefit beyond 4.5 hours (OR 1.15, 95\% CI 0.95--1.40). In absolute terms, this corresponds to roughly ten additional patients per 100 treated achieving an excellent outcome if treated within 3 hours, falling to about five per 100 between 3 and 4.5 hours. Alteplase increased the risk of fatal intracranial haemorrhage in the first week by about 2 percentage points, a risk that did not diminish with later treatment, so the balance of benefit and harm worsens as time passes. IVT can be delivered at any acute stroke unit with access to brain imaging, but it is relatively ineffective at recanalising large, proximal clots.

% MT

Mechanical thrombectomy (MT) is suitable for patients with large vessel occlusion (LVO) of the anterior circulation, who account for a minority of ischaemic strokes but a disproportionate share of death and severe disability. The HERMES collaboration pooled individual data on 1{,}287 patients from five trials and found that MT substantially reduced disability at 90 days (OR 2.49, 95\% CI 1.76--3.53), with functional independence (mRS 0--2) in 46.0\% of treated patients against 26.5\% of controls, and a number needed to treat of 2.6 for one patient to improve by at least one mRS level \cite{goyal2016}. This makes MT one of the most effective interventions in acute medicine.

As with IVT, the effect diminishes with delay. In the MR CLEAN trial, Fransen et al.\ found that the absolute benefit of MT on functional independence fell by about 6 percentage points for every hour of delay to reperfusion, from an estimated 25.9\% with reperfusion at 3 hours to 6.7\% at 6 hours \cite{fransen_time_2016}. The HERMES time-to-treatment meta-analysis confirmed the pattern across trials: each hour of delay to reperfusion was associated with worse disability (common OR 0.84, 95\% CI 0.76--0.93), and the benefit of MT over IVT was no longer statistically significant beyond 7.3 hours from onset \cite{saver2016}. Later trials showed that patients selected by advanced imaging can still benefit up to 24 hours after they were last known to be well \cite{nogueira2018,albers2018}, but these patients are a selected subgroup, and earlier treatment remains better for them too.

% Aims of project
The evidence above treats the two treatments, and the two treatment times, largely in isolation. The alteplase trials relate benefit to time to IVT in populations for whom MT was not available \cite{emberson_effect_2014}, and the MT trials relate benefit to time to reperfusion in patients, most of whom had already received alteplase, without modelling when the alteplase was given \cite{goyal2016,saver2016}. Trials comparing MT with and without prior IVT recruited patients presenting directly to MT centres, in whom the interval between the two treatments was short and varied little \cite{majoie2023,kaesmacher2024}.

In routine practice, the two times vary independently: many patients receive IVT at their local hospital and are then transferred for MT, which may follow one or several hours later. It is therefore unclear how much each treatment adds to the other at a given pair of treatment times. It is also uncertain how well trial estimates, obtained in selected patients, describe the benefit achieved in the wider population treated in routine care.
 
Large national registries offer a means of addressing these questions, but treatment in routine care is not randomised, and those treated differ systematically from those who are not. Target trial emulation, in which the analysis of observational data is designed to mirror the trial that would ideally have been run, combined with modern causal machine learning methods, allows treatment effects to be estimated from such data under explicit assumptions \cite{hernan2016}.
 
The aim of this study was to estimate the real-world effectiveness of IVT and MT, alone and in combination, in patients with ischaemic stroke and an NIHSS score above 5 (or a recorded large vessel occlusion) arriving at hospital within six hours of stroke onset, using national registry data from England and Wales. Specifically, we set out to:
 
\begin{enumerate}
    \item estimate the effect of IVT, MT, and their combination on disability at discharge, compared with no reperfusion treatment;
    \item describe how that effect depends on the time from onset to each treatment, considering time to IVT and time to MT both separately and jointly;
    \item assess the robustness of these estimates by applying three causal inference methods with differing assumptions (target trial emulation by matching, an explainable S-learner, and a causal forest), and by comparing the results with those of the randomised trials where such trials exist;
    \item identify the patient and service characteristics that most influence the benefit gained from treatment.
\end{enumerate}